\documentclass[conference]{IEEEtran}
% Joint i-SAIRAS & iSpaRo 2026. IEEEtran conference format (as used by ICRA/IROS).
% Build: pdflatex paper && bibtex paper && pdflatex paper && pdflatex paper
% Figures live in ../figs/ (generated by scripts/rung*.py and ladder_summary_figure.py).

\usepackage[T1]{fontenc}
\usepackage[utf8]{inputenc}   % allows the accented author names in biblio.bib
\usepackage{amsmath,amssymb}
\usepackage{graphicx}
\usepackage{booktabs}
\usepackage{array}
\usepackage[hidelinks]{hyperref}
\usepackage{xcolor}

\graphicspath{{../figs/}}

\begin{document}

\title{When Does Excitation Design Help Spacecraft Inertia\\ Identification? A First-Principles Study of the\\ Observability--Accuracy Link}

\author{\IEEEauthorblockN{M. El Hariry \textit{et al.}}
\IEEEauthorblockA{\textit{(affiliation)} \\ \texttt{matteo.elhariry@uni.lu}}}

\maketitle

\begin{abstract}
A recurring premise in spacecraft self-characterization is that an agent can \emph{shape} its reaction-wheel (RW) torque to enrich the observability of its inertia tensor, and that an observability functional of the excitation---e.g.\ the Fisher-information log-determinant---is therefore the right thing to optimize. We test this premise from the ground up. Starting from a clean linear-Gaussian least-squares baseline, where maximizing observability is provably identical to minimizing estimation variance, we add one realism increment at a time (measurement noise, actuator constraints, unmodeled disturbance, time-varying inertia, and finally a deployed nonlinear EKF) and measure, at each step, the rank correlation between observability and achieved inertia-estimation error. The link is a theorem at the bottom ($\rho=-0.93$), \textbf{collapses} at the first realistic increment because differentiating noisy rate measurements makes least squares biased ($\rho=-0.09$), and is \textbf{restored} only by a consistent, noise-aware estimator ($\rho=-0.87$). Once restored, excitation design becomes consequential only under actuator constraints; in loose regimes any decent broadband profile is near-optimal, but under a \emph{tight} momentum budget a closed-loop excitation that maximizes information \emph{subject to the actuator's model-validity envelope} beats the best fixed profile $10\times$ ($0.30\%$ vs.\ $3.0\%$) by avoiding the wheel saturation the fixed profile cannot. The synthesis---\emph{estimator first, observability second, excitation third, all subject to the filter's model-validity envelope}---explains why hand-tuned broadband excitation is hard to beat in benign regimes and pinpoints exactly when active sensing pays off. The errors-in-variables mechanism is classical in system identification; our contribution is to show that it specifically invalidates the Fisher-information excitation-design criterion used in the spacecraft literature unless the estimator is consistent.
\end{abstract}

\begin{IEEEkeywords}
spacecraft inertia identification, active sensing, optimal excitation, observability, Fisher information, extended Kalman filter, errors-in-variables.
\end{IEEEkeywords}

\section{Introduction}
Spacecraft inertia identification underpins attitude control, on-orbit servicing, and post-deployment reconfiguration. A natural and recurring idea is \emph{active sensing}: the spacecraft emits informative RW torques while an estimator updates its belief over the inertia tensor, and the excitation is chosen---possibly adaptively, e.g.\ by reinforcement learning---to maximize an observability functional of the resulting motion. The implicit assumptions are (A1) that observability predicts estimation accuracy, and (A2) that shaping the signal to maximize it yields a worthwhile accuracy gain.

These assumptions are rarely tested in isolation. Prior work either (i)~optimizes an open-loop excitation against a Fisher-information cost~\cite{zhai2017optimal,calafiore2001robot,swevers1997optimal}, (ii)~tunes an estimator with an observability metric~\cite{evain2024improving}, or (iii)~learns or derives a policy and reports end-to-end error~\cite{enders2021deep,elhariry2025towards}. End-to-end results entangle the excitation, the estimator, the actuator constraints, and the disturbance model, so they cannot say \emph{which} assumption holds or fails.

This paper refines our own prior study~\cite{elhariry2025towards}, which proposed actively shaping RW torque to enrich inertia observability; here we step back and ask whether that premise actually holds, isolating the conditions under which it does. We use a \textbf{ladder of controlled experiments}: each rung adds exactly one realism increment to a baseline in which the observability--accuracy relationship is analytically known, and measures how the relationship changes. Our contributions:
\begin{enumerate}
\item \textbf{A provable anchor.} In the linear-Gaussian least-squares regime, observability \emph{is} accuracy (Cram\'er--Rao with equality); we confirm it numerically and use it as ground truth.
\item \textbf{A precise failure boundary.} The link collapses at the first realistic increment---noisy rate measurements make the regression an errors-in-variables (EIV) problem and least squares becomes bias-dominated, which the Fisher information cannot see. Only a consistent estimator restores it. While EIV bias is classical~\cite{soderstrom2007errors,ljung1999system}, the consequence is specific: the Fisher-information excitation-design criterion is not a valid accuracy surrogate for spacecraft inertia ID unless paired with a consistent estimator.
\item \textbf{The locus of the design problem.} Excitation choice is irrelevant unconstrained and only becomes consequential under RW momentum/torque limits. In loose regimes any decent broadband profile is near-optimal; under a tight budget, a \emph{closed-loop in-envelope} excitation beats the best fixed profile $10\times$.
\item \textbf{A model-validity caveat for the deployed filter.} Validating against a production augmented EKF, observability is additionally gated by the filter's model-validity envelope: information-greedy excitation saturates the wheels, which the filter does not model, and the link inverts until the envelope is respected.
\end{enumerate}

\section{Related Work}
\textbf{Classical inertia identification.} Spacecraft inertia is estimated predominantly by (recursive) least squares on Euler's equation: Tanygin and Williams~\cite{tanygin1997mass} from coasting maneuvers; Bordany et al.~\cite{bordany2000inorbit} in orbit on UoSAT-12; Psiaki~\cite{psiaki2005estimation} via nested least squares from flight data; Keim et al.~\cite{keim2006constrained} as a constrained (LMI/SDP) problem enforcing physical realizability; and Norman et al.~\cite{norman2011inorbit} jointly with momentum-actuator alignment. Crucially, these obtain $\dot{\omega}$ from differentiated or integrated rate data---the operation our rung~1 shows turns the regression into an EIV problem.

\textbf{Adaptive control with embedded ID.} Ahmed et al.~\cite{ahmed1998adaptive} achieve adaptive attitude tracking with simultaneous inertia identification, using \emph{periodic commands} to render the inertia identifiable; modern concurrent-learning controllers~\cite{zhao2020finite} deliver inertia as a by-product of tracking. These assume the closed-loop task provides ``enough'' excitation.

\textbf{Optimal excitation and experiment design.} Designing an input to minimize parameter uncertainty is classical optimal experiment design~\cite{fedorov1972theory,pukelsheim2006optimal,goodwin1977dynamic}. Robotics applies it to calibration via Fourier-series excitation~\cite{swevers1997optimal} and log-det-FIM trajectory optimization~\cite{calafiore2001robot}; Zhai et al.~\cite{zhai2017optimal} transfer it to spacecraft, and Evain and Delavault~\cite{evain2024improving} use an E-optimal criterion to tune the EKF. All adopt the FIM as an accuracy surrogate---the assumption our ladder isolates.

\textbf{System identification and EIV.} That noisy regressors bias least squares, and that instrumental-variable / consistent estimators remove it, is foundational~\cite{ljung1999system,soderstrom1989system}; Söderström~\cite{soderstrom2007errors} surveys the EIV problem.

\textbf{Estimation and attitude dynamics.} Our estimators draw on the Kalman filter~\cite{kalman1960new}, nonlinear/optimal-estimation texts~\cite{simon2006optimal,crassidis2011optimal}, and spacecraft attitude/RW frameworks~\cite{markley2014fundamentals}. The disturbance-absorbing augmented state is the inertia-ID analogue of gyro-bias augmentation.

\textbf{Active SysID, learning, differentiable simulation.} Recent work makes excitation adaptive: ASID~\cite{memmel2024asid} and control-oriented active learning~\cite{lee2024active,wagenmaker2020active} drive exploration with information measures; differentiable simulation enables gradient-based identification~\cite{jatavallabhula2021gradsim}; RL handles unknown spacecraft inertia~\cite{enders2021deep,elhariry2025drl}; and Candan and Servadio~\cite{candan2026online} jointly estimate pose and full inertia of a non-cooperative target. These motivate the active-sensing premise; we test whether it pays off for cooperative inertia ID and find it does so only conditionally.

\section{Problem Formulation}
\textbf{Dynamics.} With body rate $\omega\in\mathbb{R}^3$, symmetric inertia $I$, three aligned RWs of scalar inertia $I_\mathrm{rw}$ and speeds $\Omega$, commanded wheel acceleration $u$, and an unmodeled body torque $\tau_\mathrm{ext}$,
\begin{equation}
I\dot{\omega} = \tau_\mathrm{ext} - \omega\times(I\omega + I_\mathrm{rw}\Omega) - I_\mathrm{rw} u, \qquad \dot{\Omega}=u,
\end{equation}
with RW torque/speed saturation. The six unique inertia entries are $\theta=(I_{xx},I_{yy},I_{zz},I_{xy},I_{xz},I_{yz})$. Linearizing in $\theta$ gives a $3\times6$ regressor $R(\omega,\dot\omega)$ with $R(\omega,\dot\omega)\,\theta=\tau_\mathrm{eff}$.

\textbf{Observability.} The Fisher information is $F=\sum_k R_k^\top R_k/\sigma^2$; we use D-optimality ($\log\det F$) and report A-optimality ($\mathrm{tr}\,F^{-1}$) where the theory is exact. Under disturbance we augment to estimate $\tau_\mathrm{ext}$ jointly via $[R\,|\,{-}I_3]$.

\textbf{Estimators (the ladder's controlled variable).} (i)~ordinary least squares; (ii)~instrumental variables (IV) with a time-lagged instrument, consistent under EIV; (iii)~Savitzky--Golay-smoothed LS; (iv)~recursive least squares with forgetting (for time-varying $\theta$); (v)~a 15-state augmented EKF $[\omega,I,\Omega,\tau_\mathrm{ext}]$ representative of a deployed flight estimator.

\textbf{Metric.} Frobenius relative error $\lVert\hat I-I\rVert/\lVert I\rVert$, averaged over Monte-Carlo noise draws, and its Spearman rank correlation with $\log\det F$ across a fixed family of 14 energy-matched torque profiles (single-tone, multi-tone, chirp, PRBS, random, constant).

\section{Method: The Realism Ladder}
Each rung changes exactly one factor (Table~\ref{tab:ladder}). All share one MicroSat-class satellite with asymmetric $I$, $dt=0.1$\,s, 60\,s episodes, gyro $\sigma_\omega=10^{-4}$\,rad/s.

\section{Results}
Fig.~\ref{fig:ladder} summarizes the central result: the correlation $\rho$ between observability and inertia accuracy across the ladder. Table~\ref{tab:ladder} gives the numbers.

\begin{figure}[t]
\centering
\includegraphics[width=\columnwidth]{ladder_summary.png}
\caption{The observability--accuracy correlation across the ladder. A theorem at rung~0, broken by measurement noise (1--2), restored by a consistent estimator (3+), sustained through constraints/disturbance/time-variation (4--7), and validated on the real EKF in-envelope (9) but inverted when information-greedy excitation saturates the wheels.}
\label{fig:ladder}
\end{figure}

\begin{table}[t]
\caption{The ladder at a glance. $\rho=$ Spearman(\,$\log\det F$, inertia rel-err) across the 14-profile family.}
\label{tab:ladder}
\centering
\renewcommand{\arraystretch}{1.15}
\begin{tabular}{@{}clll@{}}
\toprule
\# & increment & estimator & $\rho$ \\
\midrule
0 & clean LS (anchor)        & LS            & $-0.93$ \\
1 & gyro noise, finite-diff $\dot\omega$ & LS & $-0.09$ \\
2 & integral regression      & LS            & $-0.20$ \\
3 & (estimator change)       & IV / smoothed & $-0.87$ \\
4 & RW momentum budget       & IV            & $-0.82$ \\
5 & disturbance (augmented)  & aug.\ IV      & $-0.81$ \\
6 & optimize excitation      & aug.\ IV      & --- \\
7 & time-varying inertia     & windowed IV   & $-0.88$ \\
8 & time-varying (recursive) & RLS           & --- \\
9 & deployed EKF (in-env.)   & aug.\ EKF     & $-0.42$ \\
\bottomrule
\end{tabular}
\end{table}

\textbf{Rung 0 --- anchor.} A-optimality is exact (empirical/predicted MSE $=0.994$, $\rho(\mathrm{tr}\,F^{-1},\mathrm{MSE})=+1.00$); D-optimality ranks accuracy at $\rho=-0.93$. Under matched torque-\emph{energy}, the most observable profiles build the most angular velocity (the off-diagonals enter only through $\omega\times(I\omega)$), so the unconstrained problem is trivial.

\textbf{Rung 1 --- collapse.} Gyro noise with finite-difference $\dot\omega$ makes $R$ noisy. Least squares becomes \emph{bias}-dominated (bias/std $\approx 10\times$) and $\rho$ drops to $-0.09$: the Fisher information models only variance and is blind to the bias.

\textbf{Rung 2 --- intrinsic.} An integral (momentum-form) regression that never differentiates is $70\times$ more accurate noise-free yet still uncorrelated with observability ($\rho=-0.20$): the breakdown is intrinsic to regressing on noisy states.

\textbf{Rung 3 --- restoration.} A consistent IV estimator restores $\rho=-0.87$ and cuts the bias $20\times$ ($0.43\%\to0.02\%$); smoothing reaches $\rho=-0.86$.

\textbf{Rung 4 --- where design matters.} Under a momentum budget the trivial spin-up optimum is infeasible and the across-profile error spread explodes from ${<}2.6\%$ to $0.06\%$--$182\%$; broadband multi-axis profiles win, single-tone fail ($\rho=-0.82$).

\textbf{Rung 5 --- disturbance.} A constant $\tau_\mathrm{ext}$ inflates un-augmented error $10$--$100\times$ (median $3\%$); augmenting via $[R|{-}I_3]$ recovers it (median $0.10\%$). Joint observability of $(\theta,\tau_\mathrm{ext})$ requires rich excitation ($\rho=-0.81$).

\textbf{Rung 6 --- adaptation adds little (loose regime).} A 500-trial search for the most-observable excitation raises $\log\det F$ ($5.1\to7.4$) but improves accuracy by only $+1\%$ over the best fixed profile: the curve flattens at the noise floor.

\textbf{Rungs 7--8 --- time-varying inertia.} Under a $43\%$ ramp, a windowed tracker \emph{appears} to make the optimum regime-dependent (chirp $14\%$ vs.\ PRBS $63\%$); a recursive RLS tracker closes the gap (both $\approx3.2\%$) and improves tracking $5$--$50\times$. The apparent adaptive headroom was an estimator artifact.

\textbf{Rung 9 --- deployed EKF.} The production augmented EKF, from a biased $85\%$-of-true initial inertia under disturbance, reproduces the story in-envelope ($0.3$--$0.8\%$). But the open-loop FIM \emph{mispredicts} accuracy when excitation saturates the wheels---the filter models $\dot\Omega=u$ with no saturation---and the correlation inverts to $+0.34$, returning to $-0.42$ once saturation is removed.

\textbf{Rung 10 --- closed-loop in-envelope excitation.} Motivated by rung~9, we run a greedy controller that picks the next torque to maximize information in the EKF's least-observed inertia direction \emph{subject to a hard wheel-momentum envelope} (Fig.~\ref{fig:cl}). In a loose budget it ties the best fixed profile. Under a \emph{tight} momentum budget it wins decisively: $0.30\%$ rel-err with $0\%$ saturation, versus $3.0\%$ for the best fixed broadband profile (which saturates $46\%$ of the time), and versus $7.6\%$ for the same controller \emph{without} the envelope (which spins up and saturates $91\%$ of the time). Closing the loop helps precisely when fixed profiles cannot respect the actuator envelope.

\begin{figure}[t]
\centering
\includegraphics[width=\columnwidth]{rung10_closed_loop.png}
\caption{Rung 10, tight momentum budget. Closed-loop in-envelope excitation beats the best fixed broadband profile $10\times$ by avoiding wheel saturation; removing the envelope constraint reproduces the rung-9 saturation failure.}
\label{fig:cl}
\end{figure}

\section{Discussion}
\textbf{Estimator first.} Across every regime, switching to a consistent/augmented/recursive estimator improved accuracy by 1--2 orders of magnitude---far more than any excitation difference. The estimator is the dominant lever; observability becomes meaningful only once it is correct.

\textbf{Observability is valid but bounded.} Given a noise-aware estimator, $\log\det F$ ranks excitation accuracy ($\rho\approx-0.8$ to $-0.9$). But for the deployed nonlinear filter it must be maximized \emph{subject to the model-validity envelope} (no wheel saturation, bounded $\omega$).

\textbf{When adaptation pays.} In benign regimes the in-envelope optimum is flat and pinned at the noise floor, so explicit/learned excitation adds little---consistent with the empirical observation that hand-tuned broadband matches learned policies. But under tight actuator constraints, where fixed open-loop profiles violate the momentum envelope, closed-loop in-envelope excitation is decisively better (rung~10). This delimits exactly where active sensing is worth its complexity.

\textbf{Limitations.} Single satellite and disturbance class; constant-$\tau_\mathrm{ext}$ augmentation; simulation only. Closed-loop excitation against the estimated dynamics under richer disturbance models is the natural next study.

\section{Conclusion}
Excitation design for spacecraft inertia identification helps---but second-order, and conditionally. The estimator (consistent, augmented, recursive) is the load-bearing element; observability is a valid excitation criterion only with a noise-aware estimator and only within the deployed filter's model-validity envelope; decent broadband excitation is near-optimal in benign regimes, while closed-loop in-envelope excitation earns its keep under tight actuator constraints. We offer the ladder as a reproducible template for separating estimator, excitation, and observability effects that end-to-end studies conflate.

\bibliographystyle{IEEEtran}
\bibliography{biblio}

\end{document}
